\documentclass[11pt]{article}

\usepackage[utf8]{inputenc}
\usepackage[T1]{fontenc}
\usepackage[brazil]{babel}
\usepackage{amsmath, amssymb}
\usepackage{graphicx}
\usepackage{geometry}
\usepackage{fancyhdr}
\usepackage{enumitem}
%\usepackage{physics}
\usepackage{hyperref}
\usepackage{multicol}
\usepackage{xcolor}
\usepackage{siunitx}
\usepackage{comment}
\usepackage{amsthm}

\theoremstyle{definition}
\newtheorem{obs}{Observação}
\newtheorem{definicao}{Definição}
\newtheorem{teorema}{Teorema}
\newtheorem{exemplo}{Exemplo}

\geometry{a4paper, margin=2.5cm}

\pagestyle{fancy}
\fancyhf{}
\rhead{Campos conservativos}
\lhead{Caio Nicácio\,\href{https://orcid.org/0009-0009-2877-6170}{\raisebox{-0.2\height}{
\includegraphics[height=1.8ex]{../../orcid.pdf}}}}
\cfoot{\thepage}

\setlength{\parskip}{0.5em}
\setlength{\parindent}{0pt}


\begin{document}
	
	\begin{center}
		\Large\textbf{CAMPOS CONSERVATIVOS} \\
		\normalsize
		Data: \today
	\end{center}
	
	\hrule
	\vspace{1em}
	
	\tableofcontents
	\newpage
	
	\section{Teorema fundamental do cálculo para integrais de linha}
	\label{sec:tfc}
	
	Seja $f:[a,b]\to \mathbb{R}$, $f\in C^1$, uma função diferenciável em $]a,b[$.
	O primeiro teorema fundamental do cálculo estabelece que
	\begin{equation*}
		\int_{a}^{b} f(x)\, dx=F(b)-F(a),
	\end{equation*}
	onde $F$ é uma primitiva de $f$.
	
	Assim, a integral no intervalo $[a,b]$ é obtida avaliando um das primitivas nos extremos do intervalo.
	
	\begin{definicao} \textbf{(Conjunto conexo)}
		Um conjunto aberto $\Omega\subset \mathbb{R}^n$ é simplesmente conexo, se existe uma curva suave $\gamma:I \to \Omega$ por partes ligado a dois pontos quaisquer de $\Omega$. Logo, $\forall \boldsymbol{a}, \boldsymbol{b}  \in \Omega$, $\exists \gamma:[a,b] \to \Omega \subset \mathbb{R}^n$, tal que $\boldsymbol{a}=\gamma(a)$ e $\boldsymbol{b}=\gamma(b)$.
	\end{definicao}
	\begin{definicao} \textbf{(Conjunto desconexo)}
		Um conjunto aberto $\Omega \subset \mathbb{R}^n$ é dito desconexo se $\Omega$ é a união de dois ou mais conjuntos abertos disjuntos e não vazios, ou seja,
		\begin{equation*}
			\Omega=\bigcup_{i \in \Gamma}A_i
		\end{equation*}
		com $A_i\cap A_j = \emptyset$, $i \neq j$, em que $\Gamma$ denota um conjunto de índices e $A_i\neq \emptyset$ é aberto $\forall i \in \Gamma$.
	\end{definicao}
	\begin{definicao}
		A integral de linha\footnote{Para informações complementares, clique \href{https://github.com/caio-nicacio/studyNotes/blob/main/calculus-topics/line-integrals/line-integrals.pdf}{aqui}.} $\int_{\gamma} \vec{\boldsymbol{F}}\cdot d\gamma$, $\vec{\boldsymbol{F}}:\Omega \subset \mathbb{R}^n\to \mathbb{R}^n$ contínuo no conjunto aberto conexo $\Omega$ é independente do caminho se
		\begin{equation*}
			\int_{\gamma_1} \vec{\boldsymbol{F}}\cdot d\gamma_1=\int_{\gamma_2} \vec{\boldsymbol{F}}\cdot d\gamma_2=\dots=\int_{\gamma_k} \vec{\boldsymbol{F}}\cdot d\gamma_k,\quad i=1,2,\dots, k,
		\end{equation*}
		$\forall \gamma_i:I_i$, $\gamma_i$ fechada.
	\end{definicao}
	\begin{teorema}
		A integral de linha $\int_{\gamma}\vec{\boldsymbol{F}}\cdot d\gamma$ de um campo vetorial contínuo $\vec{\boldsymbol{F}}:\Omega\subset \mathbb{R}^n\to \mathbb{R}^n$ definido sobre um conjunto simplesmente conexo $\Omega$ é independente do caminho se, e somente se, $\int_{\gamma}\vec{\boldsymbol{F}}\cdot d\gamma=0$, $\gamma:[a,b]\to \Omega$, $\gamma$ curva fechada.
	\end{teorema}
	\begin{teorema} \label{thm:tfc2} \textbf{(Segundo teorema fundamental do cálculo)}
		Seja $\varphi:\Omega\to \mathbb{R}$ uma função potencial de $\vec{\boldsymbol{F}}$, tal que $\nabla \varphi$ é contínuo em um conjunto aberto simplesmente conexo $\Omega\subset \mathbb{R}^n$. Dada uma curva suave por partes $\gamma:[a,b]\to \Omega$, tem-se que
		\begin{equation*}
			\int_{\gamma}\nabla \varphi\cdot d\gamma = \varphi(\gamma(b))-\varphi(\gamma(b)).
		\end{equation*}
	\end{teorema}
	Seja $\gamma:[a,b]\to \Omega$ uma curva suave e seja $\varphi:\Omega\to \mathbb{R}$ função potencial de $\vec{\boldsymbol{F}}$. Define-se a função composta $g(t)=\varphi(\gamma(t))$, $\forall t\in [a,b]$.
	
	Pela regra da cadeia, $g'(t)=\nabla \varphi(\gamma(t))\cdot \gamma'(t)$. Além disso, visto que $\gamma$ é suave e $\nabla \varphi$ contínuo, $g'(t)$ é continua.
	
	Aplicando o Teorema \ref{thm:tfc2} à função  $g$, conclui-se que
	\begin{gather*}
		\int_{\gamma}\nabla \varphi\cdot d\gamma=\int_{a}^{b} \nabla \varphi(\gamma(t))\cdot \gamma'(t)\, dt=\int_{a}^{b} g'(t)\, dt\\
		\int_{\gamma}\nabla \varphi\cdot d\gamma=g(b)-g(a)\\
		\int_{\gamma}\nabla \varphi\cdot d\gamma=\varphi(\gamma(b))-\varphi(\gamma(a)).
	\end{gather*}
	Na condição de $\gamma:[a,b]\to \Omega$ ser uma curva suave por partes, particiona-se $[a,b]$ em $k$ subintervalos $[t_{i-1},t_i]$, $i=1,2,\dots,k$, e aplica-se o resultado anterior a cada segmento.
	\begin{align*}
		\int_{\gamma} \nabla \varphi\cdot d\gamma&=\int_{a}^{b} \nabla \varphi(\gamma(t)\cdot \gamma'(t)\, dt)\\
		&=\sum_{i=1}^{n}\int_{t_{i-1}}^{t_i}\nabla \varphi(\gamma(t))\cdot \gamma'(t)\, dt\\
		&=\sum_{i=1}^{n}[\varphi(\gamma(t_i))-\varphi(\gamma(t_{i-1}))]\\
		&=\varphi(\gamma(t_k))-\varphi(\gamma(t_0))\\
		\int_{\gamma} \nabla \varphi\cdot d\gamma&=\varphi(\gamma(b))-\varphi(\gamma(a)).
	\end{align*}

	\section{Campos conservativos}
	
	Um campo vetorial $\vec{\boldsymbol{F}}:\Omega \subset \mathbb{R}^n \to \mathbb{R}^n$ denomina-se conservativo se existe um campo escalar diferenciável $\varphi:\Omega \to \mathbb{R}$ tal que
	\begin{gather*}
		\nabla \varphi=\vec{\boldsymbol{F}}\\
		\frac{\partial \varphi}{\partial x}\,\mathbf{i}+\frac{\partial \varphi}{\partial y}\,\mathbf{j}+\frac{\partial \varphi}{\partial z}\,\mathbf{k}=P(x,y,z)\,\mathbf{i}+Q(x,y,z)\,\mathbf{j}+R(x,y,z)\mathbf{k}.
	\end{gather*}

	\subsection{Função potencial}
	\begin{quote}
		Função escalar cujo gradiente gera um campo vetorial, o qual formaliza a proposição de que o campo vetorial pode ser descrito como a variação espacial de uma grandeza escalar.
	\end{quote}
	Seja $\Omega \subset \mathbb{R}^n$ um domínio aberto e seja $\vec{\boldsymbol{F}}:\Omega\subset \mathbb{R}^n \to \mathbb{R}^n$ um campo vetorial contínuo. Diz-se que $\varphi:\Omega\to \mathbb{R}$ é uma função potencial de $\vec{\boldsymbol{F}}$ se $\nabla \varphi=\vec{\boldsymbol{F}}$. Isto é,
	\begin{equation*}
		\vec{\boldsymbol{F}}(x_1,x_2,\dots,x_n)=\left(\frac{\partial \varphi}{\partial x_1},\frac{\partial \varphi}{\partial x_2},\dots,\frac{\partial \varphi}{\partial x_n}\right).
	\end{equation*}
	Nesse caso, $\vec{\boldsymbol{F}}$ é denominado campo conservativo.
	
	Consequentemente, se $\vec{\boldsymbol{F}}=\nabla \varphi$, então a integral de linha entre dois pontos $A$ e $B$ é independente do caminho. Logo,
	\begin{equation*}
		\int_{\gamma}\vec{\boldsymbol{F}}\cdot d\vec{\boldsymbol{r}}=\varphi(B)-\varphi(A),
	\end{equation*}
	$\forall\gamma$ suave ligando $A$ e $B$.
	
	Em domínios simplesmente conexos, $\vec{\boldsymbol{F}}=\nabla \varphi$ se, e somente se, $\operatorname{rot} \vec{\boldsymbol{F}}=0$.
	
	\begin{teorema}
		Seja $\vec{\boldsymbol{F}}:\Omega \subset \mathbb{R}^n\to \mathbb{R}^n$, $n=2,3$, um campo vetorial de classe $C^1$ no aberto $\Omega$. Uma condição necessária (mas não suficiente) para $\vec{\boldsymbol{F}}$ ser considerado conservativo é
		\begin{equation*}
			\operatorname{rot}\vec{\boldsymbol{F}}=0\, \text{em}\, \Omega.
		\end{equation*} 
	\end{teorema}
	Seja $\vec{\boldsymbol{F}}:\Omega\to \mathbb{R}^3$ um campo conservativo, $\exists\varphi:\Omega\to \mathbb{R}$ tal que $\nabla\varphi$ em $\Omega$, e
	\begin{equation*}
		\frac{\partial \varphi}{\partial x}=P(x,y,z),\, \frac{\partial \varphi}{\partial y}=Q(x,y,z),\, \frac{\partial \varphi}{\partial z}=R(x,y,z).
	\end{equation*}
	Visto que $\vec{\boldsymbol{F}}\in C^1$, resulta que $\varphi\in C^2$. Logo,
	\begin{equation*}
		\frac{\partial^2 \varphi}{\partial x \partial y}=\frac{\partial^2 \varphi}{\partial y \partial x}\quad \text{(Teorema de Schwarz)}.
	\end{equation*}
	Segue que
	\begin{equation*}
		\frac{\partial P}{\partial y}=\frac{\partial Q}{\partial x}\quad \longrightarrow\quad \frac{\partial P}{\partial y}-\frac{\partial Q}{\partial x}=0.
	\end{equation*}
	Analogamente,
	\begin{equation*}
		\frac{\partial P}{\partial y}=\frac{\partial R}{\partial x}\quad \text{e}\quad \frac{\partial Q}{\partial z}=\frac{\partial R}{\partial y},
	\end{equation*}
	em $\Omega$. Logo, $\operatorname{rot}\vec{\boldsymbol{F}}=0$, $\vec{\boldsymbol{F}}:\Omega \subset \mathbb{R}^3$.
	\begin{obs}
		A condição é topológica, não dependendo apenas do campo $\vec{\boldsymbol{F}}$, mas do conjunto $\Omega$, no qual deve ser simplesmente conexo para que $\oint \vec{\boldsymbol{F}}\cdot d\gamma=0$.
	\end{obs}
	\begin{comment}
		\begin{figure}[h!]
			\begin{minipage}[b]{0.48\textwidth}
				\centering
				\includegraphics[width=\textwidth]{figure/estado-geral-tensao-1.png}
				\\[-1ex]\small (a)
				\label{fig:estado-geral-1}
			\end{minipage}
			\hfill
			\begin{minipage}[b]{0.48\textwidth}
				\centering
				\includegraphics[width=\textwidth]{figure/estado-geral-tensao-2.png}
				\\[-1ex]\small (b)
				\label{fig:estado-geral-2}
			\end{minipage}
			\centering
			%\captionsetup{width=16cm}
			\caption{\centering Representação do estado geral de tensões em um ponto $Q$: (a) no sistema de coordenadas original e (b) no sistema resultante de uma rotação em relação ao sistema inicial.}
			\label{fig:estado-geral}
		\end{figure}
	
		\begin{figure}[h!]
			\centering
			\includegraphics[width=5.5cm]{figure/parametrizacao.jpeg}
			%\captionsetup{width=16cm}
			\caption{\centering Parametrização arbitrária $\boldsymbol{\vec{r}}(t)$ cuja imagem é a curva $\gamma(t):[a, b]$.}
			\label{fig:parametrizacao}
		\end{figure}
	\end{comment}
	
	%\bibliography{pos-textuais/referencias.bib}
	
\end{document}
